% ============================================================
%  Book chapter 19: Dates and Times
% ============================================================
\chapter{Dates and Times}

\textbf{What this chapter gives you:} the four types you need from \code{datetime}
(\code{date}, \code{time}, \code{datetime}, \code{timedelta}), the difference between a
\code{str}, a \code{datetime} and a float of seconds, and a few rules that stop time zone bugs.
In AI work you meet this in log timestamps, API \code{created} fields, run durations and dated
dataset files.

\tikzset{c19n/.style={draw=#1,fill=#1!10,rounded corners=3pt,minimum height=8mm,inner sep=4pt,
    font=\ttfamily\small},
  c19t/.style={draw=#1,fill=#1!10,minimum height=7mm,inner sep=3pt,font=\ttfamily\small}}

% ------------------------------------------------------------
\section{Why dates and times are hard}

The same day can be written many ways, and some are ambiguous.

\begin{center}
\begin{tabularx}{\linewidth}{@{}l X@{}}
\toprule
\textbf{Text} & \textbf{Problem} \\
\midrule
\code{21 Jul 1987} & month name depends on the language \\
\code{7/21/1987} & only the US order; 21 cannot be a month \\
\code{03/04/2025} & 4 March or 3 April? Depends on the country \\
\code{2025-03-04} & ISO 8601: year, month, day. Unambiguous, and it sorts as text \\
\bottomrule
\end{tabularx}
\end{center}

Times have a worse problem: time zones follow politics, not geometry, and daylight saving time
(DST) moves the clock on different days in different countries. One moment in time has many
local clock readings.

\begin{center}
\begin{tikzpicture}[>=Stealth]
  \node[c19n=cblue,align=center] (u) at (0,0) {one instant\\[-1pt]{\scriptsize\sffamily 2025-03-14 09:00 UTC}};
  \node[c19n=cgreen] (a) at (6.5,1.1) {14:30 \ \ India (+05:30)};
  \node[c19n=cgold]  (b) at (6.5,0) {09:00 \ \ UTC (+00:00)};
  \node[c19n=cpurple] (c) at (6.5,-1.1) {05:00 \ \ New York (-04:00)};
  \draw[->,thick,cgray] (u.east) -- (a.west);
  \draw[->,thick,cgray] (u.east) -- (b.west);
  \draw[->,thick,cgray] (u.east) -- (c.west);
\end{tikzpicture}
\end{center}
% MD-BEGIN
% ```mermaid
% flowchart LR
%     U["one instant<br/>2025-03-14 09:00 UTC"] --> A["14:30 India (+05:30)"]
%     U --> B["09:00 UTC (+00:00)"]
%     U --> C["05:00 New York (-04:00)"]
% ```
% MD-END

DST makes it worse. In spring one local hour does not exist. In autumn one local hour happens
twice. Here is New York in 2025:

\begin{center}
\begin{tikzpicture}[>=Stealth]
  \node[font=\sffamily\scriptsize,anchor=east] at (-0.2,1.0) {9 March (spring)};
  \node[c19t=cblue] at (0.5,1.0) {01:30};
  \node[c19t=cred,dashed,fill=cred!5,minimum width=26mm,text=cred] at (3.0,1.0) {02:xx missing};
  \node[c19t=cblue] at (5.5,1.0) {03:00};
  \draw[->,thick,cgray] (1.2,1.0) -- (1.7,1.0);
  \draw[->,thick,cgray] (4.4,1.0) -- (4.9,1.0);
  \node[font=\sffamily\scriptsize,anchor=east] at (-0.2,-0.2) {2 November (autumn)};
  \node[c19t=cblue] at (0.5,-0.2) {01:30};
  \node[c19t=cgold,minimum width=20mm] at (2.6,-0.2) {01:30 again};
  \node[c19t=cblue] at (5.1,-0.2) {02:00};
  \draw[->,thick,cgray] (1.2,-0.2) -- (1.5,-0.2);
  \draw[->,thick,cgray] (3.7,-0.2) -- (4.5,-0.2);
  \node[font=\sffamily\tiny,text=cgray] at (0.5,-0.7) {EDT};
  \node[font=\sffamily\tiny,text=cgray] at (2.6,-0.7) {EST (clock set back)};
\end{tikzpicture}
\end{center}
% MD-BEGIN
% ```text
% 9 March 2025 (spring forward):   01:30 -> [02:xx does not exist] -> 03:00
% 2 Nov 2025 (fall back):          01:30 (EDT) -> 01:30 (EST, again!) -> 02:00
% ```
% MD-END

\begin{gamerule}[: store UTC, show local]
Keep every stored or computed timestamp in UTC (or as an epoch number). Convert to a local zone
only at the edge, when a human reads it. Servers should run on UTC. Local wall-clock times with
DST create hours that do not exist and hours that happen twice.
\end{gamerule}

Python spreads dates over several modules: \code{datetime} (the main one), \code{time} (epoch
seconds and clocks) and \code{calendar} (calendar facts). They overlap, which is confusing at first.

% ------------------------------------------------------------
\section{Leap Year}

A year has about 365.2422 days, so calendars add a day now and then. The rule has three
layers, checked from the most special case outward.

\begin{center}
\begin{tikzpicture}[c19q/.style={draw=cblue,fill=cblue!8,rounded corners=3pt,minimum height=8mm,
    minimum width=30mm,font=\ttfamily\footnotesize},
  c19r/.style={draw=#1,fill=#1!12,rounded corners=3pt,minimum height=7mm,minimum width=18mm,font=\sffamily\small\bfseries},
  >=Stealth]
  \node[c19q] (q1) at (0,0) {year \% 400 == 0};
  \node[c19q] (q2) at (4.6,0) {year \% 100 == 0};
  \node[c19q] (q3) at (9.2,0) {year \% 4 == 0};
  \node[c19r=cgreen] (r1) at (0,-1.7) {leap};
  \node[c19r=cred] (r2) at (4.6,-1.7) {not leap};
  \node[c19r=cgreen] (r3) at (9.2,-1.7) {leap};
  \node[c19r=cred] (r4) at (13.0,0) {not leap};
  \draw[->,thick] (q1) -- node[left,font=\sffamily\scriptsize]{yes} (r1);
  \draw[->,thick] (q1) -- node[above,font=\sffamily\scriptsize]{no} (q2);
  \draw[->,thick] (q2) -- node[left,font=\sffamily\scriptsize]{yes} (r2);
  \draw[->,thick] (q2) -- node[above,font=\sffamily\scriptsize]{no} (q3);
  \draw[->,thick] (q3) -- node[left,font=\sffamily\scriptsize]{yes} (r3);
  \draw[->,thick] (q3) -- node[above,font=\sffamily\scriptsize]{no} (r4);
\end{tikzpicture}
\end{center}
% MD-BEGIN
% ```mermaid
% flowchart LR
%     Q1{"year % 400 == 0"} -- yes --> L1["leap"]
%     Q1 -- no --> Q2{"year % 100 == 0"}
%     Q2 -- yes --> N1["not leap"]
%     Q2 -- no --> Q3{"year % 4 == 0"}
%     Q3 -- yes --> L2["leap"]
%     Q3 -- no --> N2["not leap"]
% ```
% MD-END

\begin{pycode}
import calendar

print([(y, calendar.isleap(y)) for y in (1900, 1996, 2000, 2023, 2024)])
print(calendar.monthrange(2024, 2)[1], calendar.monthrange(2023, 2)[1])
print(calendar.day_name[0], calendar.month_name[2])
\end{pycode}
\begin{rawcode}
[(1900, False), (1996, True), (2000, True), (2023, False), (2024, True)]
29 28
Monday February
\end{rawcode}

\code{monthrange(year, month)} returns the weekday of the first day and the number of days. Do not
write your own leap rule. Use \code{calendar.isleap}, or better, let \code{datetime} do the date
arithmetic so leap days are handled for you.

% ------------------------------------------------------------
\section{The datetime Module}

\code{datetime} has four main classes.

\begin{center}
\begin{tabularx}{\linewidth}{@{}l X l@{}}
\toprule
\textbf{Class} & \textbf{Holds} & \textbf{Looks like} \\
\midrule
\code{date} & year, month, day & \code{2025-03-14} \\
\code{time} & hour, minute, second, microsecond & \code{09:30:00} \\
\code{datetime} & a date and a time together & \code{2025-03-14T09:30:00} \\
\code{timedelta} & a length of time & \code{1 day, 3:00:00} \\
\bottomrule
\end{tabularx}
\end{center}

\begin{center}
\begin{tikzpicture}[>=Stealth]
  \node[c19n=cgreen] (d) at (0,1.0) {date};
  \node[c19n=cgreen] (t) at (0,-0.6) {time};
  \node[c19n=cblue] (dt) at (5.2,0.2) {datetime};
  \node[c19n=cgold] (td) at (10.6,0.2) {timedelta};
  \draw[->,thick,cgray] (d.east) -- node[above,sloped,font=\ttfamily\scriptsize]{combine} ([yshift=2mm]dt.west);
  \draw[->,thick,cgray] (t.east) -- node[below,sloped,font=\ttfamily\scriptsize]{combine} ([yshift=-2mm]dt.west);
  \draw[->,thick,cgray,bend left=25] (dt.north east) to node[above,font=\ttfamily\scriptsize]{dt2 - dt1} (td.north west);
  \draw[->,thick,cgray,bend left=25] (td.south west) to node[below,font=\ttfamily\scriptsize]{dt + td} (dt.south east);
\end{tikzpicture}
\end{center}
% MD-BEGIN
% ```mermaid
% flowchart LR
%     D["date"] -- combine --> DT["datetime"]
%     T["time"] -- combine --> DT
%     DT -- "dt2 - dt1" --> TD["timedelta"]
%     TD -- "dt + td" --> DT
%     DT -- ".date() / .time()" --> D
% ```
% MD-END

\subsection{date and timedelta}

\begin{pycode}
from datetime import date, datetime, timedelta, timezone
from datetime import time as dtime   # avoid a clash with the time module

d = date(2025, 3, 14)
print(d, repr(d))
print(d.year, d.month, d.day, d.weekday(), d.strftime("%A"))
print(d + timedelta(days=30), d - timedelta(weeks=2))
gap = date(2025, 12, 25) - d
print(gap, type(gap), gap.days)
print(date.min, date.max)
two = timedelta(days=2)
print(date(2024, 2, 28) + two, date(2025, 2, 28) + two)
\end{pycode}
\begin{rawcode}
2025-03-14 datetime.date(2025, 3, 14)
2025 3 14 4 Friday
2025-04-13 2025-02-28
286 days, 0:00:00 <class 'datetime.timedelta'> 286
0001-01-01 9999-12-31
2024-03-01 2025-03-02
\end{rawcode}

\code{weekday()} counts Monday as 0, so Friday is 4. \code{isoweekday()} counts Monday as 1. The
leap day in 2024 shows up with no extra code.

\begin{typecheck}
\code{print(d)} shows \code{2025-03-14}, but \code{d} is a \code{date} object, not a string.
\code{repr(d)} shows its real nature. Subtracting two dates gives a \code{timedelta}, not an
\code{int}: use \code{.days} or \code{.total\_seconds()} to get a number.
\end{typecheck}

A \code{timedelta} stores days, seconds and microseconds. Careful with \code{.seconds}: it is only
the seconds part after whole days, not the full length.

\begin{pycode}
a = datetime(2025, 3, 14, 12, 0, 0)
b = datetime(2025, 3, 14, 12, 0, 7, 250000)
td = b - a
print(td, td.total_seconds(), td.seconds, td.microseconds)
print(timedelta(days=1, hours=3), timedelta(minutes=90).total_seconds())
print(timedelta(days=1.5), timedelta(hours=3) / timedelta(hours=2))
print(timedelta(hours=25) // timedelta(hours=2), timedelta(hours=1) * 3)
print(timedelta(seconds=-5))
\end{pycode}
\begin{rawcode}
0:00:07.250000 7.25 7 250000
1 day, 3:00:00 5400.0
1 day, 12:00:00 1.5
12 3:00:00
-1 day, 23:59:55
\end{rawcode}

\begin{trap}
Negative timedeltas print strangely: five seconds backward is \code{-1 day, 23:59:55}. It is
correct, just stored as days plus a positive remainder. Use \code{total\_seconds()} when you want
\code{-5.0}. Also, \code{.seconds} of a 2-day delta is \code{0}, not 172800.
\end{trap}

\begin{pycode}
for bad in (lambda: a - d, lambda: d + d):
    try:
        bad()
    except TypeError as e:
        msg = str(e).replace("unsupported operand type(s) for ", "")
        print("TypeError:", msg)
\end{pycode}
\begin{rawcode}
TypeError: -: 'datetime.datetime' and 'datetime.date'
TypeError: +: 'datetime.date' and 'datetime.date'
\end{rawcode}

\begin{center}
\begin{tabularx}{\linewidth}{@{}l l X@{}}
\toprule
\textbf{Expression} & \textbf{Result type} & \textbf{Note} \\
\midrule
\code{date - date} & \code{timedelta} & also \code{datetime - datetime} \\
\code{date + timedelta} & \code{date} & \code{datetime + timedelta} gives \code{datetime} \\
\code{timedelta * int} & \code{timedelta} & \\
\code{timedelta / timedelta} & \code{float} & \code{//} gives an \code{int} \\
\code{date + date} & \code{TypeError} & adding two points in time makes no sense \\
\code{datetime - date} & \code{TypeError} & do not mix the two types \\
\bottomrule
\end{tabularx}
\end{center}

\subsection{time and datetime}

\begin{pycode}
t = dtime(9, 30)
print(repr(t), t.hour, t.isoformat())
dt = datetime(2025, 1, 2, 3, 4, 5, 6)   # y, m, d, h, min, s, microsecond
print(repr(dt))
print(dt.isoformat())
print(dt.isoformat(sep=" ", timespec="seconds"))
full = datetime.combine(d, t)
print(full, repr(full.date()), full.time())
print(full.replace(hour=0, minute=0), full.replace(year=2026))
print(type(datetime.now()), type(date.today()))
\end{pycode}
\begin{rawcode}
datetime.time(9, 30) 9 09:30:00
datetime.datetime(2025, 1, 2, 3, 4, 5, 6)
2025-01-02T03:04:05.000006
2025-01-02 03:04:05
2025-03-14 09:30:00 datetime.date(2025, 3, 14) 09:30:00
2025-03-14 00:00:00 2026-03-14 09:30:00
<class 'datetime.datetime'> <class 'datetime.date'>
\end{rawcode}

Objects of these classes are \emph{immutable}. \code{replace()} returns a new object. The
\code{now()} and \code{today()} calls give the current moment, so their results change each time
you run them; I show only their types above.

\begin{trap}
\code{datetime.now()} returns local time \emph{without} any time zone attached. Two machines in
different zones produce different values for the same instant, and Python cannot tell which is
which. See the section on aware datetimes below.
\end{trap}

\subsection{Naive versus aware}

A \code{datetime} is \emph{naive} when it has no \code{tzinfo}, and \emph{aware} when it carries
a UTC offset. Naive values mean ``some wall clock, somewhere''. Aware values mean one exact
instant.

\begin{pycode}
naive = datetime(2025, 3, 14, 9, 0)
aware = datetime(2025, 3, 14, 9, 0, tzinfo=timezone.utc)
print(naive.tzinfo, aware.tzinfo, aware.isoformat())
try:
    naive < aware
except TypeError as e:
    print("TypeError:", e)
ist = timezone(timedelta(hours=5, minutes=30), "IST")
local = aware.astimezone(ist)
print(local.isoformat(), local == aware)
print(datetime.now(timezone.utc).tzinfo)
\end{pycode}
\begin{rawcode}
None UTC 2025-03-14T09:00:00+00:00
TypeError: can't compare offset-naive and offset-aware datetimes
2025-03-14T14:30:00+05:30 True
UTC
\end{rawcode}

\code{astimezone} changes how the instant is \emph{shown}, not which instant it is, so
\code{local == aware} is \code{True}. For real zones with DST rules, use the standard
\code{zoneinfo} module. On Windows it needs the \code{tzdata} package (\code{pip install tzdata}); I
could not run that here.

\begin{pycode}
# needs: tzdata on Windows (pip install tzdata); not run in these notes
from zoneinfo import ZoneInfo

ny = ZoneInfo("America/New_York")
instant = datetime(2025, 3, 14, 9, 0, tzinfo=timezone.utc)
print(instant.astimezone(ny))      # the same instant on a New York wall clock
\end{pycode}

\begin{boss}[: naive versus aware]
Most date bugs in real systems are naive/aware bugs. A naive \code{datetime} is just numbers; it
does not say which zone it is in. An aware one is an exact instant. You cannot compare or subtract
the two kinds (Python raises \code{TypeError}, which is the helpful outcome). The safe pattern:
create aware UTC values (\code{datetime.now(timezone.utc)}), store them as ISO strings with an
offset or as epoch numbers, and convert with \code{astimezone} only for display.
\end{boss}

% ------------------------------------------------------------
\section{The time Module}

The \code{time} module is \emph{not} the same as the \code{datetime.time} class. It works with
\emph{epoch seconds}: the number of seconds since midnight UTC on 1 January 1970. This is the
simplest way to pass a moment between systems, such as JavaScript, databases or API fields
(\code{"created": 1700000000}).

\begin{center}
\begin{tikzpicture}[>=Stealth]
  \draw[->,thick,cgray] (0,0) -- (12.5,0);
  \foreach \x/\lab/\val in {0.3/1970-01-01 00:00:00 UTC/0, 5.2/2001-09-09 01:46:40 UTC/1{,}000{,}000{,}000, 9.6/2023-11-14 22:13:20 UTC/1{,}700{,}000{,}000} {
    \draw[thick,cblue] (\x,-0.12) -- (\x,0.12);
    \node[font=\ttfamily\small,above=3pt] at (\x,0.1) {\val};
    \node[font=\sffamily\scriptsize,text=cgray,below=4pt,align=center] at (\x,-0.1) {\lab};
  }
  \node[font=\sffamily\scriptsize,text=cgray,anchor=west] at (11.2,0.4) {seconds};
\end{tikzpicture}
\end{center}
% MD-BEGIN
% ```text
% 0                        1,000,000,000              1,700,000,000        ---> seconds
% |------------------------|--------------------------|
% 1970-01-01 00:00:00 UTC   2001-09-09 01:46:40 UTC     2023-11-14 22:13:20 UTC
% ```
% MD-END

\begin{pycode}
import calendar, time

st = time.gmtime(1_700_000_000)          # UTC, as a struct_time
print(type(st).__name__, tuple(st))
print(time.asctime(st))
print(st.tm_year, st[0], st.tm_yday, st.tm_wday)
print(time.strftime("%Y-%m-%d %H:%M:%S", st))
print(calendar.timegm(st))               # back to epoch seconds (UTC)
print(type(time.time()), type(time.perf_counter()))
\end{pycode}
\begin{rawcode}
struct_time (2023, 11, 14, 22, 13, 20, 1, 318, 0)
Tue Nov 14 22:13:20 2023
2023 2023 318 1
2023-11-14 22:13:20
1700000000
<class 'float'> <class 'float'>
\end{rawcode}

\code{time.time()} returns a \code{float} of epoch seconds. \code{localtime(secs)} gives a
\code{struct\_time} in your zone, \code{gmtime(secs)} in UTC. A \code{struct\_time} is a named tuple
of nine fields: year, month, day, hour, minute, second, weekday, day of year, DST flag. Weekday
counts from 0 (Monday), day of year from 1.

\begin{center}
\begin{tabularx}{\linewidth}{@{}l l X@{}}
\toprule
\textbf{Function} & \textbf{Returns} & \textbf{Use it for} \\
\midrule
\code{time.time()} & \code{float}, epoch seconds & timestamps to store or send \\
\code{time.gmtime(s)} / \code{localtime(s)} & \code{struct\_time} & breaking seconds into parts \\
\code{time.mktime(st)} & \code{float} & \code{struct\_time} in local zone back to seconds \\
\code{calendar.timegm(st)} & \code{int} & \code{struct\_time} in UTC back to seconds \\
\code{time.perf\_counter()} & \code{float} & measuring how long code takes \\
\code{time.monotonic()} & \code{float} & timeouts that must never go backward \\
\code{time.sleep(s)} & \code{None} & pause for \code{s} seconds \\
\bottomrule
\end{tabularx}
\end{center}

\begin{gamerule}
To time code, use \code{time.perf\_counter()}, not \code{time.time()}. The wall clock can jump
(clock sync, DST), while \code{perf\_counter} only moves forward and has high resolution. The
absolute value is meaningless; subtract two readings.
\end{gamerule}

\begin{pycode}
t0 = time.perf_counter()
time.sleep(0.05)
elapsed = time.perf_counter() - t0
print(0.04 < elapsed < 0.2)
\end{pycode}
\begin{rawcode}
True
\end{rawcode}

% ------------------------------------------------------------
\section{Read and Write Dates and Times}

\subsection{ISO 8601: the default choice}

\code{isoformat()} writes it, \code{fromisoformat()} reads it. In Python 3.11 and later it
accepts a trailing \code{Z} too. Use ISO for logs, JSON, JSONL eval files and file names: it is
unambiguous and sorts correctly as plain text.

\begin{pycode}
print(date.fromisoformat("2025-03-14") == d)
print(sorted(["2025-03-14", "2025-01-09", "2024-12-31"]))
print(datetime.fromisoformat("2025-02-14T15:05:09+00:00"))
print(repr(datetime.fromisoformat("2025-02-14T15:05:09Z")))
\end{pycode}
\begin{rawcode}
True
['2024-12-31', '2025-01-09', '2025-03-14']
2025-02-14 15:05:09+00:00
datetime.datetime(2025, 2, 14, 15, 5, 9, tzinfo=datetime.timezone.utc)
\end{rawcode}

\subsection{strftime: object to string}

\code{strftime} (``string from time'') formats a date or time with percent codes.
\code{strptime} (``string parsed to time'') goes the other way with the same codes.

\begin{center}
\begin{tikzpicture}[>=Stealth]
  \foreach \x/\f/\r in {0/{\%Y}/2025, 1.6/{-}/{-}, 2.4/{\%m}/02, 3.8/{-}/{-}, 4.6/{\%d}/14, 6.0/{ }/{ }, 6.8/{\%H}/15, 8.2/{:}/{:}, 9.0/{\%M}/05} {
    \node[font=\ttfamily\small] at (\x,0.9) {\f};
    \node[font=\ttfamily\small] at (\x,-0.2) {\r};
    \draw[->,cgray] (\x,0.6) -- (\x,0.15);
  }
  \node[font=\sffamily\scriptsize,anchor=west,text=cgray] at (10.0,0.9) {format string};
  \node[font=\sffamily\scriptsize,anchor=west,text=cgray] at (10.0,-0.2) {result};
\end{tikzpicture}
\end{center}
% MD-BEGIN
% ```text
% format string:  %Y  -  %m  -  %d  space  %H  :  %M
%                  |      |      |           |      |
% result       :  2025 - 02  -  14   space   15  :  05      ->  "2025-02-14 15:05"
% ```
% MD-END

\begin{center}
\begin{tabularx}{\linewidth}{@{}l l l X@{}}
\toprule
\textbf{Code} & \textbf{Meaning} & \textbf{Example} & \textbf{Note} \\
\midrule
\code{\%Y} \ \code{\%y} & year 4 / 2 digits & \code{2025} \ \code{25} & \\
\code{\%m} & month & \code{02} & zero-padded \\
\code{\%d} & day of month & \code{14} & zero-padded \\
\code{\%H} \ \code{\%I} \ \code{\%p} & hour 24 / 12, AM or PM & \code{15} \ \code{03} \ \code{PM} & \\
\code{\%M} \ \code{\%S} & minute, second & \code{05} \ \code{09} & \\
\code{\%f} & microseconds & \code{123456} & 6 digits \\
\code{\%b} \ \code{\%B} & month name short / long & \code{Feb} \ \code{February} & locale dependent \\
\code{\%a} \ \code{\%A} & weekday short / long & \code{Fri} \ \code{Friday} & locale dependent \\
\code{\%j} & day of year & \code{045} & \\
\code{\%z} \ \code{\%Z} & UTC offset, zone name & \code{+0000} \ \code{UTC} & aware values only \\
\bottomrule
\end{tabularx}
\end{center}

\begin{pycode}
dt = datetime(2025, 2, 14, 15, 5, 9)
fmt = "It's %A, %B %d, %Y, %I:%M:%S%p"
print(dt.strftime(fmt))
print(dt.strftime("%Y%m%d-%H%M%S"))
print(date(2025, 2, 14).strftime("%d/%m/%Y %H:%M"))
print(dtime(10, 35).strftime("%H:%M"), dtime(10, 35).strftime("%Y-%m-%d"))
\end{pycode}
\begin{rawcode}
It's Friday, February 14, 2025, 03:05:09PM
20250214-150509
14/02/2025 00:00
10:35 1900-01-01
\end{rawcode}

\begin{trap}
A \code{date} has no time, so \code{\%H:\%M} prints midnight (\code{00:00}). A \code{time} has no
date, so \code{\%Y-\%m-\%d} prints the meaningless \code{1900-01-01}. Format only the parts the
object really has.
\end{trap}

\subsection{strptime: string to object}

\code{datetime.strptime(text, fmt)} returns a \code{datetime}. It is not a regex: the characters
that are not percent codes (dashes, spaces, slashes) must match exactly. It also checks that the
date exists.

\begin{pycode}
p = datetime.strptime("2025-02-14 15:05:09", "%Y-%m-%d %H:%M:%S")
print(repr(p), p == dt, type(p))
print(datetime.strptime("14 Feb 2025", "%d %b %Y").date())
print(datetime.strptime("03/04/2025", "%d/%m/%Y").date())
print(datetime.strptime("03/04/2025", "%m/%d/%Y").date())
for text in ("2025 02 14", "2025-02-30"):
    try:
        datetime.strptime(text, "%Y-%m-%d")
    except ValueError as e:
        print("ValueError:", e)
\end{pycode}
\begin{rawcode}
datetime.datetime(2025, 2, 14, 15, 5, 9) True <class 'datetime.datetime'>
2025-02-14
2025-04-03
2025-03-04
ValueError: time data '2025 02 14' does not match format '%Y-%m-%d'
ValueError: day is out of range for month
\end{rawcode}

\begin{trap}
The same text \code{03/04/2025} gives 3 April or 4 March depending on the format you choose, and
Python will not warn you. Never guess the order. Ask the data source, or demand ISO. Also,
\code{strptime} returns a \emph{naive} datetime unless the text contains an offset parsed with
\code{\%z}.
\end{trap}

\subsection{Locale and names}

Month and weekday names come from the operating system's locale. \code{locale.setlocale(locale.LC\_TIME, "fr\_FR")}
switches them. That works only if the OS has that locale installed (names vary, such as
\code{fr\_FR.UTF-8} on Linux). It changes process-wide state, so avoid it in servers. For storing
and exchanging data, use numeric or ISO formats, which do not depend on language.

% ------------------------------------------------------------
\section{All the Conversions}

You end up converting among a handful of forms all the time. The \code{datetime} object is the hub.

\begin{center}
\begin{tikzpicture}[>=Stealth]
  \node[c19n=cblue] (dt) at (5.5,0) {datetime};
  \node[c19n=cgreen] (iso) at (0,1.7) {str (ISO)};
  \node[c19n=cgold] (fmt) at (0,-1.7) {str (custom)};
  \node[c19n=cpurple] (ep) at (11,1.7) {epoch float};
  \node[c19n=cteal] (dd) at (11,-1.7) {date + time};
  \draw[<->,thick,cgray] (iso) -- node[above,sloped,font=\ttfamily\scriptsize]{isoformat / fromisoformat} (dt);
  \draw[<->,thick,cgray] (fmt) -- node[below,sloped,font=\ttfamily\scriptsize]{strftime / strptime} (dt);
  \draw[<->,thick,cgray] (ep) -- node[above,sloped,font=\ttfamily\scriptsize]{timestamp / fromtimestamp} (dt);
  \draw[<->,thick,cgray] (dd) -- node[below,sloped,font=\ttfamily\scriptsize]{.date() .time() / combine} (dt);
\end{tikzpicture}
\end{center}
% MD-BEGIN
% ```mermaid
% flowchart LR
%     ISO["str (ISO)"] <-- "isoformat / fromisoformat" --> DT["datetime"]
%     FMT["str (custom)"] <-- "strftime / strptime" --> DT
%     EP["epoch float"] <-- "timestamp / fromtimestamp" --> DT
%     DD["date + time"] <-- ".date() .time() / combine" --> DT
% ```
% MD-END

\begin{pycode}
a = datetime(2025, 3, 14, 9, 0, tzinfo=timezone.utc)
print(a.timestamp())
print(datetime.fromtimestamp(1_700_000_000, tz=timezone.utc).isoformat())
print(a.strftime("%Y-%m-%dT%H:%M:%SZ"))
print(a.timetuple().tm_yday)
\end{pycode}
\begin{rawcode}
1741942800.0
2023-11-14T22:13:20+00:00
2025-03-14T09:00:00Z
73
\end{rawcode}

\begin{trap}
\code{datetime.fromtimestamp(secs)} \emph{without} \code{tz=} gives a naive datetime in your local
zone. For API timestamps, always pass \code{tz=timezone.utc}. Likewise \code{timestamp()} on a naive
datetime assumes local time, which is a silent source of off-by-hours errors.
\end{trap}

A short AI-flavored example: duration of a run from ISO log timestamps.

\begin{pycode}
lines = [
    "2025-03-14T09:00:01Z start run=a",
    "2025-03-14T09:00:03.500Z first token run=a",
    "2025-03-14T09:02:11Z done run=a",
]
stamps = [datetime.fromisoformat(l.split()[0]) for l in lines]
print([str(y - x) for x, y in zip(stamps, stamps[1:])])
print((stamps[-1] - stamps[0]).total_seconds())
print(date(1990, 5, 17) + timedelta(days=10000))
\end{pycode}
\begin{rawcode}
['0:00:02.500000', '0:02:07.500000']
130.0
2017-10-02
\end{rawcode}

The last line answers a typical ``when am I 10{,}000 days old'' question for someone born on
17 May 1990.

% ------------------------------------------------------------
\section{Alternative Modules}

If the standard library feels clumsy, these packages exist. None is needed for this series.

\begin{center}
\begin{tabularx}{\linewidth}{@{}l X@{}}
\toprule
\textbf{Package} & \textbf{What it is good for} \\
\midrule
\code{python-dateutil} & parses almost any date text; relative dates such as ``next month'' \\
\code{arrow} & one friendly object for create, shift, format \\
\code{whenever} & newer, fast library (written in Rust) with clear naive/aware types \\
\code{iso8601} & fills gaps in ISO 8601 parsing \\
\code{astropy.time}, \code{TAI64} & astronomy-size dates, atomic time \\
\bottomrule
\end{tabularx}
\end{center}

\begin{gamerule}
Before you add a date library, try \code{datetime} with \code{timezone.utc} and
\code{fromisoformat}. It covers most log and API work. Use \code{dateutil.parser} only for messy
human-written dates you cannot control.
\end{gamerule}

% ------------------------------------------------------------
\section{Review and what comes next}

You can now tell a date, a time, a datetime and a timedelta apart, convert between ISO strings,
custom strings, epoch numbers and objects, and you know to keep UTC at the core. The next chapter
opens files: reading, writing and moving data on disk.

\begin{practice}
\begin{enumerate}[nosep,leftmargin=*]
  \item Which of the years 1800, 2000, 2023 and 2100 are leap years? Check with code.
  \item Create the date of your birthday, print its ISO form and its weekday name, and compute the
        date on which you turn 10{,}000 days old.
  \item Why does \code{date(2025, 3, 1) - date(2025, 2, 1)} have type \code{timedelta} and not
        \code{int}? Which attribute gives a number of days?
  \item Parse \code{"07/08/2025"} two times, once as day-first and once as month-first. Print both
        results in ISO format.
  \item Write one line that turns the API field \code{1700000000} into an ISO string in UTC.
  \item What happens if you subtract \code{datetime.now()} from \code{datetime.now(timezone.utc)}?
        Why is that error a good thing?
\end{enumerate}
\end{practice}

\begin{achievement}{Time Keeper}
\begin{itemize}[nosep,leftmargin=*]
  \item Choose the right type among \code{date}, \code{time}, \code{datetime} and \code{timedelta}.
  \item Parse and format timestamps with ISO, \code{strptime} and \code{strftime}.
  \item Convert between epoch seconds, aware datetimes and strings.
  \item Explain naive versus aware datetimes and why you store UTC.
  \item Time a piece of code correctly with \code{perf\_counter}.
\end{itemize}
\end{achievement}
